\documentclass[10pt]{article}

\usepackage[margin=0.75in]{geometry}
\usepackage{amsmath,amsthm,amssymb,bm}
\usepackage{xcolor}
\usepackage{cancel}
\usepackage{graphicx}
\usepackage{changepage}
\usepackage{circuitikz}
\usepackage{pgfplots}
\usepackage{minted}
\usepackage{physics}
\usepackage{hyperref}
\usepackage{siunitx}
\usepackage{fontspec}
\usepackage{relsize}
\usepackage{subfig}
\usepackage{todonotes}
\usepackage{multicol, multirow, booktabs}
\usepackage[breakable]{tcolorbox}
\usepackage[inline]{enumitem}

\theoremstyle{definition}
\newtheorem{problem}{Problem}
\newtheorem{soln}{Solution}

\pgfplotsset{compat=newest}
\usetikzlibrary{lindenmayersystems}
\usetikzlibrary{arrows}
\usetikzlibrary{calc}
\usetikzlibrary{positioning, fit}
\usetikzlibrary{3d, perspective}

\definecolor{incolor}{HTML}{303F9F}
\definecolor{outcolor}{HTML}{D84315}
\definecolor{cellborder}{HTML}{CFCFCF}
\definecolor{cellbackground}{HTML}{F7F7F7}
\newcommand{\ui}{\hat{i}}
\newcommand{\uj}{\hat{j}}
\newcommand{\uk}{\hat{k}}
\newcommand{\ux}{\hat{\mathbf{x}}}
\newcommand{\uy}{\hat{\mathbf{y}}}
\newcommand{\uz}{\hat{\mathbf{z}}}
\newcommand{\uv}[1]{\hat{\bv{{#1}}}}
\newcommand{\pr}[1]{{#1^\prime}}
\newcommand{\pd}[2][]{{\frac{\partial^{#1}}{\partial {{#2}^{#1}}}}}
\newcommand{\dif}[2][]{{\frac{d^{#1}}{d{{#2}^{#1}}}}}
\pgfdeclarelayer{background}  
\pgfsetlayers{background,main}
\AtBeginDocument{\RenewCommandCopy\qty\SI}

\makeatletter
\newcommand{\boxspacing}{\kern\kvtcb@left@rule\kern\kvtcb@boxsep}
\makeatother
\newcommand{\prompt}[4]{
    \ttfamily\llap{{\color{#2}[#3]:\hspace{3pt}#4}}\vspace{-\baselineskip}
}

\newcommand{\thevenin}[2]{
  \begin{center}
    \begin{circuitikz} \draw
      (0,0) -- (2,0) to[battery1, l_=$V_{Th}\eq#1$] (2,2) 
      to[resistor, l_=$R_{Th}\eq#2$] (0,2)
      ;
      \draw [o-] (-.07,2.079);
      \draw [o-] (-.07,0.079);
    \end{circuitikz}
  \end{center}
}

\newcommand{\norton}[2]{
  \begin{center}
    \begin{circuitikz} \draw
      (0,0) -- (3,0) to[american current source, l_=$I_{N}\eq#1$] (3,2) -- (0,2) (2,0)
      to[resistor, l=$R_{N}\eq#2$] (2,2)
      ;
      \draw [o-] (-.07,2.079);
      \draw [o-] (-.07,0.079);
    \end{circuitikz}
  \end{center}
}

\DeclareMathOperator{\Div}{div}
\DeclareMathOperator{\Curl}{curl}
\DeclareMathOperator{\sgn}{sgn}

\DeclareSIUnit\foot{ft}

\newcommand{\highlight}[1]{\colorbox{yellow}{$\displaystyle #1$}}

\newcommand{\ti}[1]{\widetilde{#1}}

\newfontface{\Kaufmann}{Kaufmann}
\DeclareTextFontCommand{\kf}{\Kaufmann}
\newcommand{\scriptr}{\fontsize{12pt}{12pt}\kf{r}}

\newfontface{\KaufmannB}{Kaufmann Bold}
\DeclareTextFontCommand{\kfb}{\KaufmannB}
\newcommand{\bscriptr}{\fontsize{12pt}{12pt}\kfb{r}}
\newcommand{\justif}[2]{&{#1}&\text{#2}}
\newcommand{\bv}[1]{\bm{\mathrm{#1}}}


\title{Physics 3200Y: Lab III Prelab}
\author{Jeremy Favro (0805980) \\ Trent University, Peterborough, ON, Canada}
\date{\today}

\begin{document}
\maketitle

% PROBLEM 1
\begin{problem}
For a circular coil of $N$ turns that is mounted on a shaft that rotates about the axis that is along
the shaft: Estimate the minimum number of turns $N$ required (and its equivalent length of wire $L$) for a $2.5,
  3.0,3.5,\dots,\qty{9.0}{\centi\meter}$ diameter coil $d$ to produce $> \qty{0.02}{\newton\meter}$ instantaneous torque $\tau$ driven by $\qty{1.5}{\ampere}$ pulsed current
$I$ (regardless of AWG) under a uniform $B$-field (along the coil plane) from an AlNiCo magnet. Assume that the
force, wire/current direction, and $B$-field are all perpendicular to each other at the time of pulse/contact.
\end{problem}
\begin{soln}
  As given in the manual,
  $$\bv{F}=I\bv{L}\cross\bv{B}$$
  and the cross product of perpendicular vectors is
  $$F=ILB.$$
  We can get the torque then as
  $\tau=I\frac{\pi d^2B}{4}$ and the torque felt by $N$ loops as $N$ times the torque felt by one,
  $$\tau=IN\frac{\pi d^2B}{4}.$$
  So if want a specific torque we need
  $$N=\frac{4\tau}{I\pi d^2B}=\frac{8}{I\pi d^2 B}$$
  where $d$ is in meters.
  Hence we obtain,

  \begin{table}[ht!]
    \centering%
    \begin{tabular}{ll}
      \toprule
      $d\,\unit{\centi\meter}$ & $N$    \\
      \midrule
      2.5                      & 181.08 \\
      3                        & 125.75 \\
      3.5                      & 92.39  \\
      4                        & 70.74  \\
      4.5                      & 55.89  \\
      5                        & 45.27  \\
      5.5                      & 37.41  \\
      6                        & 31.44  \\
      6.5                      & 26.79  \\
      7                        & 23.1   \\
      7.5                      & 20.12  \\
      8                        & 17.68  \\
      8.5                      & 15.66  \\
      9                        & 13.97  \\
      \bottomrule
    \end{tabular}
  \end{table}
\end{soln}

% PROBLEM 2
\begin{problem}
or a circular coil of $\pr{N}$ turns as a voice coil that is $\qty{1.6}{\centi\meter}$ in diameter $\pr{d}$: Estimate the minimum
length of wire $\pr{L}$ (and its equivalent number of turns $\pr{N}$) required for the AWG sizes $30,32,\dots, 38$ to build a
$\qty{2}{\watt}$ loudspeaker using the audio amplifier as the voltage source ($\qty{1.5}{\volt}$ RMS). Reference the parameters/tables
necessary for your calculated estimates.
\end{problem}
\begin{soln}
  We know that $P=V^2/R$. Using the hyperphysics AWG resistance table \url{http://hyperphysics.phy-astr.gsu.edu/hbase/Tables/wirega.html} and a feet-to-meters conversion factor of $3.281$ we obtain that, for a coil of
  length $L$ and resistance per foot of $K$, $R= 3.281KL$ so
  $$\frac{V^2}{3.281 KP}=L$$
  and we know that
  $$L=CN=N\pi d\implies \frac{L}{\pi d}=N=\frac{V^2}{3.281 KP\pi d}.$$

  \begin{table}[ht!]
    \centering%
    \begin{tabular}{lll}
      \toprule
      AWG & $\unit{\ohm\per\foot}$ & N     \\
      \midrule
      30  & 0.1032                 & 66.1  \\
      32  & 0.1641                 & 41.57 \\
      34  & 0.2609                 & 26.15 \\
      36  & 0.4148                 & 16.45 \\
      38  & 0.6596                 & 10.34 \\
      \bottomrule
    \end{tabular}
  \end{table}



\end{soln}
\end{document}